\section{Experimental setup}

The main measure is real wall-clock speedup over the target model alone, measured on the same machine. The other metrics explain where that speedup comes from.

\begin{table}[h]
\centering
\caption{Metrics.}
\label{tab:metrics}
\begin{tabularx}{\linewidth}{@{}lX@{}}
\toprule
Metric & Meaning \\
\midrule
Speedup & Tokens/sec of a method $\div$ tokens/sec of the target alone (live, same PC) \\
Mean accepted length & Average tokens gained per target call ($n+1$) \\
Simulated reward & Average episode reward in the simulator, on train and test prompts \\
Training cost & Episodes and wall time each algorithm needs to converge \\
Simulation-vs-live gap & Whether the simulator ranks methods in the same order as the live benchmark \\
\bottomrule
\end{tabularx}
\end{table}

\paragraph{Hypotheses.}
\begin{enumerate}
  \item All four RL agents will match or beat the best fixed-$k$ baseline, because they can adapt draft length to the text.
  \item Value Iteration will be the strongest tabular method, since it plans with an estimated model, and Q-learning will approach it with enough samples.
  \item DQN and REINFORCE will gain from the richer continuous state (entropy, acceptance rate), with REINFORCE showing noisier training.
  \item The learned policies will look like a threshold on $q$, matching the theory.
  \item Speedups will be larger on the SmolLM2 pair (12$\times$ size gap) than on the Qwen pair (3$\times$ gap).
\end{enumerate}

\todo{hardware details of the benchmark PC, software versions (PyTorch, transformers, CUDA), training hyperparameters.}
